\chapter{Podsumowanie}
\label{chap:podsumowanie}
Niniejsza praca miała na celu stworzenie projektu aplikacji ułatwiającej współpracę przy grupowych projektach studenckich oraz pozyskiwanie potencjalnych inwestorów.

\section{Rezultat pracy}
Rezultatem pracy jest prototyp platformy wspierającej studentów w~prowadzeniu wspólnych projektów programistycznych. 
System łączy funkcje zarządzania projektami i~członkami zespołu, definiowania ról i~uprawnień, planowania pracy w~sprintach, obsługi repozytoriów Git oraz publikowania ogłoszeń naboru i~przyjmowania aplikacji. W~odróżnieniu od narzędzi skupionych na pojedynczym obszarze, zaprojektowana aplikacja integruje te funkcje w~jednym systemie.

\section{Wnioski}
Podczas implementacji projektu aplikacji autor poznał dokładniej podejście projektowe \textit{Domain-Driven Design}. Podejście to nie zostało wdrożone w~pełni, natomiast wykorzystano kilka jego głównych elementów, które pozwoliły na lepsze odzwierciedlenie domeny problemu. Pierwotnie planowano zastosować tę metodę ortodoksyjnie, lecz w~trakcie prac zauważono, że nadmiernie skomplikuje to proces tworzenia oprogramowania i~utrudni późniejsze utrzymanie projektu. W~podejściu tym sprawdziły się agregaty, encje oraz zdarzenia domenowe, które uzupełniają model anemiczny. Nie znaleziono natomiast odpowiedniego zastosowania dla obiektów wartości (\textit{Value Objects}), które w~tym przypadku okazały się zbędne. Zrezygnowano również z~rozbudowanych, bogatych modeli domenowych, gdyż agregaty w~zupełności wystarczyły do spójnego odwzorowania domeny i~przechowywania logiki biznesowej.

Elementem, który bardzo dobrze sprawdził się w~tym projekcie, było wykorzystanie architektury warstwowej z~użyciem wzorca \textit{CQRS}. Podział na komendy i~zapytania pozwolił na czystsze odseparowanie warstwy prezentacji od warstwy aplikacji. Implementacja kontrolerów w~API została znacząco uproszczona i~stała się mniej zależna od usług w~warstwie aplikacji. Uzyskano również większą czytelność kodu oraz łatwiejszą utrzymywalność projektu, dzięki rezygnacji z~tworzenia obszernych interfejsów usług. Skorzystanie z~wzorca \textit{CQRS} pozwoliło także na uproszczenie implementacji zdarzeń domenowych, które zostały utworzone z~użyciem wbudowanego mechanizmu powiadomień we wzorcu \textit{Mediator}. W~trakcie realizacji projektu nie zauważono istotnych wad tego podejścia.

Podczas tworzenia aplikacji największe wyzwanie stanowiła poprawna obsługa usługi repozytoriów kodu źródłowego, co wymagało od autora dogłębnego zrozumienia działania systemu kontroli wersji \textit{Git} oraz zapoznania się z~jego dokumentacją i~API. Była to jedna z~kluczowych funkcjonalności systemu, dlatego jej implementacja pochłonęła najwięcej czasu. Jednocześnie moduł ten może stanowić w~przyszłości pewne ograniczenie aplikacji --- miejsce na dysku serwera może stosunkowo szybko się wyczerpać i~spowodować błędy w~działaniu systemu. Przypadek ten nie został jeszcze zabezpieczony, jednak problem został zidentyfikowany i~w~razie dalszego rozwoju oprogramowania zostanie odpowiednio obsłużony.

\section{Dalsze kierunki rozwoju} 
Przed wdrożeniem produkcyjnym konieczne byłoby dalsze rozwinięcie systemu, w~szczególności zapewnienie skalowalnego przechowywania repozytoriów oraz dodanie wyszukiwania i~tagowania projektów. 
Kolejnym ważnym krokiem byłoby opracowanie wbudowanego systemu komunikacji pomiędzy inwestorami a~członkami projektu, tak aby obie strony nie były zmuszone do korzystania z~zewnętrznych kanałów, takich jak telefony komórkowe czy poczta elektroniczna. 
Istotny byłby również rozwój modułu repozytoriów w~kierunku umożliwienia użytkownikom tworzenia nowych gałęzi, wystawiania tzw. \emph{Pull Requests} oraz przeprowadzania procesu \emph{Code Review}. 
W~przyszłych pracach należałoby także ustalić priorytety funkcjonalności oraz zweryfikować wymagania dotyczące wydajności i~bezpieczeństwa.